% ============================================================
\chapter{Modul 7: Dateien und Daten}
% ============================================================

\section{Lektion 20: Dateien lesen}

\begin{lstlisting}
import json, csv

# --- JSON-Konfiguration lesen ---
def lade_config(pfad="roboter.json"):
    with open(pfad, "r", encoding="utf-8") as f:
        return json.load(f)

# Beispiel-JSON-Datei (roboter.json):
# {"name": "MobileRobot-1", "max_v": 0.5, "min_abstand": 0.2}

# --- CSV-Sensorlog lesen ---
def lade_sensorlog(pfad="log.csv"):
    messwerte = []
    with open(pfad, newline="", encoding="utf-8") as f:
        reader = csv.DictReader(f)
        for zeile in reader:
            messwerte.append({
                "zeit":    float(zeile["zeit_s"]),
                "abstand": float(zeile["abstand_m"]),
                "akku":    int(zeile["akku_prozent"])
            })
    return messwerte

# Textdatei lesen
with open("wegpunkte.txt", "r") as f:
    for zeile in f:
        x, y = zeile.strip().split(",")
        print(f"Wegpunkt: ({x}, {y})")
\end{lstlisting}

\begin{merksatz}
\texttt{with open(...) as f:} ist der sichere Weg. Die Datei wird automatisch geschlossen -- auch wenn ein Fehler auftritt. Ohne \texttt{with}: Datei bleibt offen, Datenverlust möglich.
\end{merksatz}

\section{Lektion 21: Dateien schreiben -- Sensorlogs}

\begin{lstlisting}
import csv, json, time

# --- CSV-Sensorlog schreiben ---
def starte_log(pfad="sensor_log.csv"):
    f = open(pfad, "w", newline="", encoding="utf-8")
    writer = csv.DictWriter(f,
        fieldnames=["zeit_s", "abstand_m", "akku_prozent", "modus"])
    writer.writeheader()
    return f, writer

def schreibe_messung(writer, abstand, akku, modus):
    writer.writerow({
        "zeit_s":       round(time.time(), 3),
        "abstand_m":    abstand,
        "akku_prozent": akku,
        "modus":        modus
    })

# --- JSON-Konfiguration schreiben ---
def speichere_config(config, pfad="config.json"):
    with open(pfad, "w", encoding="utf-8") as f:
        json.dump(config, f, indent=2, ensure_ascii=False)
    print(f"Konfiguration gespeichert: {pfad}")

konfig = {"name": "MobileRobot-1", "max_v": 0.5, "min_abstand": 0.20}
speichere_config(konfig)
\end{lstlisting}

\begin{praxis}
Robotersysteme loggen immer Daten: Sensorwerte, Positionen, Fehlermeldungen. Diese Logs werden später zur Analyse und Fehlersuche verwendet. CSV ist das universelle Format -- jede Tabellenkalkulation kann es öffnen.
\end{praxis}

\begin{uebung}[title=Projektaufgabe Modul 7: Datenlogger]
Schreibe einen vollständigen Datenlogger: Simuliere 20 Messzyklen, schreibe jeden Zyklus in eine CSV-Datei (Zeit, Abstand, Akku, X-Position, Y-Position). Danach: Datei lesen, Minimum/Maximum/Durchschnitt berechnen, Zusammenfassung ausgeben.
\end{uebung}
